% Lapisan lokalisasi untuk makro pembaca Fremlin.
% Muat setelah volwp.aux; tidak mengubah byte sumber otoritatif.

% Beri pemutus baris sedikit ruang darurat untuk kata majemuk Bahasa Indonesia
% dan rumus sebaris yang lebih panjang; ini hanya mereflow tata letak A4.
\emergencystretch=3em

\long\def\proof#1{\ifwithproofs{\medskip

     \noindent{\bf Bukti} #1}\else\fi}
\def\Hint#1{({\it Petunjuk\/}: #1)}

% Istilah lokal bagi fungsi yang melestarikan ukuran melalui prapeta.
\def\imp{pelestari ukuran melalui prapeta}

\def\Notesheader#1{\immediate\write0{#1 Catatan dan komentar}
     \wheader{#1 Catatan}{16}{8}{5}{48pt}
     \noindent{\bf #1 Catatan dan komentar}}
\def\vNotesheader#1#2{\immediate\write0{#2 Catatan dan komentar}
     \wheader{#2 Catatan}{16}{8}{5}{#1}
     \noindent{\bf #2 Catatan dan komentar}}

% Jangan perluas judul bermatematika di dalam \write; judul tercetak tetap utuh.
\def\leader#1#2{\immediate\write0{#1}
     \wheader{#1}{14}{7}{5}{24pt}
     {\bf #1 #2}}
\def\vleader#1#2#3{\immediate\write0{#2}
     \wheader{#2}{14}{7}{5}{#1}
     {\bf #2 #3}}

\def\newchapter#1{\gdef\topparagraph{}
   \gdef\bottomparagraph{Bab\ #1 {\it pengantar}}
   \gdef\newparagraph{Bab\ #1 {\it pengantar}}
   \footmarkcount=0
   \gdef\sectionname{Pendahuluan}
   \gdef\headlinesectionname{Pendahuluan}
   \centerline{\bf Bab #1}
   \medskip
   \centerline{\bf \chaptername}
   \medskip
   }
\def\newsection#1{\gdef\bottomparagraph{\S#1 {\it pengantar}}
   \noindent{\bf #1 \sectionname}
   \smallskip
   \let\headlinesectionname=\sectionname}
\def\newvolume#1{\gdef\topparagraph{}
   \gdef\bottomparagraph{Jilid\ #1 {\it pengantar}}
   \Loadfont{\twelvebf=cmbx12}
   \def\chaptername{Jilid #1}
   \def\sectionname{\volumename}
   \def\headlinesectionname{\volumename}
   \centerline{\twelvebf Jilid #1:  \volumename}
   \bigskip
   }

\def\leftfootline{\ifwithproofs{\eightsmc Fondasi Teori Ukuran\hfill}
    \else{\eightsmc Fondasi Teori Ukuran \eightit (versi ringkas)\hfill}
    \fi}
